











\begin{define}{Greatest Common Divisor}
  For integers $a,b$, not both zero, the \emph{greatest common divisor} $d=\gcd(a,b)$ is the largest positive integer dividing both $a$ and $b$.
\end{define}

% --- Euclidean Algorithm Section ---

\begin{theorem}[Euclidean Algorithm]
Let $a, b \in \mathbb{Z}$, not both zero. The Euclidean Algorithm is a procedure to compute $d = \gcd(a, b)$ by repeated division:
\begin{align*}
  a &= bq_1 + r_1, & 0 \leq r_1 < |b| \\
  b &= r_1 q_2 + r_2, & 0 \leq r_2 < r_1 \\
  r_1 &= r_2 q_3 + r_3, & 0 \leq r_3 < r_2 \\
  &\vdots \\
  r_{k-2} &= r_{k-1} q_k + r_k, & 0 \leq r_k < r_{k-1} \\
  r_{k-1} &= r_k q_{k+1} + 0
\end{align*}
The last nonzero remainder $r_k$ is $\gcd(a, b)$.
\end{theorem}

\begin{example}
Compute $\gcd(252, 105)$ using the Euclidean Algorithm:
\begin{align*}
  252 &= 105 \times 2 + 42 \\
  105 &= 42 \times 2 + 21 \\
  42 &= 21 \times 2 + 0
\end{align*}
The last nonzero remainder is $21$, so $\gcd(252, 105) = 21$.
\end{example}

\begin{theorem}[Extended Euclidean Algorithm]
The Extended Euclidean Algorithm not only computes $d = \gcd(a, b)$, but also finds integers $u, v$ such that $au + bv = d$ (a Bézout identity).
It works by back-substituting through the steps of the Euclidean Algorithm.
\end{theorem}

\begin{proof}
Suppose the Euclidean Algorithm for $a, b$ produces:
\begin{align*}
  a &= bq_1 + r_1 \\
  b &= r_1 q_2 + r_2 \\
  &\vdots \\
  r_{k-2} &= r_{k-1} q_k + r_k \\
  r_{k-1} &= r_k q_{k+1} + 0
\end{align*}
Then $\gcd(a, b) = r_k$. By expressing each remainder as a linear combination of $a$ and $b$ (working backwards), we can write $r_k = au + bv$ for some $u, v \in \mathbb{Z}$.
\end{proof}

\begin{example}
Find integers $u, v$ such that $252u + 105v = 21$.
From the previous example:
\begin{align*}
  252 &= 105 \times 2 + 42 \\[-1ex]
  105 &= 42 \times 2 + 21 \\[-1ex]
  42 &= 21 \times 2 + 0
\end{align*}
Work backwards:
\begin{align*}
  21 &= 105 - 42 \times 2 \\
     &= 105 - 2(252 - 2 \times 105) \\
     &= 105 - 2 \times 252 + 4 \times 105 \\
     &= -2 \times 252 + 5 \times 105
\end{align*}
So $u = -2$, $v = 5$ and $252(-2) + 105(5) = 21$.
\end{example}

% --- End Euclidean Algorithm Section ---

\begin{define}{Linear Diophantine Equation}
  A \emph{linear Diophantine equation} in two variables is an equation of the form
  \[
    a x + b y = c,
  \]
  where $a,b,c\in\mathbb Z$ and $x,y$ are integer unknowns.
\end{define}


\begin{theorem}[Existence and General Solution]
Let $a,b,c\in\mathbb Z$, not both $a,b$ zero, and set $d=\gcd(a,b)$.  Then:
\begin{enumerate}[label=\textup{(\arabic*)}]
  \item The equation
    \[
      a x + b y = c
    \]
    has an integer solution $(x,y)$ if and only if $d\mid c$.
  \item If $d\mid c$, then a particular solution $(x_0,y_0)$ can be found by
    \begin{itemize}[nosep]
      \item using the Extended Euclidean Algorithm to find $(u,v)$ such that
        \[a u + b v = d,\]
      \item and setting
        \[(x_0,y_0) = \bigl(u\,\tfrac{c}{d},\;v\,\tfrac{c}{d}\bigr).\]
    \end{itemize}
  \item All integer solutions are given by
    \[
      x = x_0 + \frac{b}{d}\,t,
      \quad
      y = y_0 - \frac{a}{d}\,t,
      \quad t\in\mathbb Z.
    \]
\end{enumerate}
\end{theorem}

\begin{proof}
\textbf{(1) Necessity.}  If $x,y\in\mathbb Z$ satisfy $a x + b y = c$, then any common divisor of $a,b$ divides the left side, hence divides $c$.  In particular $d=\gcd(a,b)\mid c$.\\

\noindent\textbf{(1) Sufficiency.}  By the Extended Euclidean Algorithm there exist $u,v\in\mathbb Z$ with
\[a u + b v = d.\]
Since $d\mid c$, write $c=d\,k$ and multiply to get
\[a(u k)+b(v k)=c,\]
so $(x_0,y_0)=(u k, v k)$ solves $a x+b y=c$.\\

\noindent\textbf{(2)} Follows directly by multiplying the Bézout combination by $c/d$.\\

\noindent\textbf{(3)} If $(x_0,y_0)$ is one solution, then for any $t\in\mathbb Z$,
\begin{align*}
  a\bigl(x_0+\tfrac{b}{d}t\bigr) + b\bigl(y_0-\tfrac{a}{d}t\bigr)
  &= (a x_0 + b y_0) + t\,(a\tfrac{b}{d}-b\tfrac{a}{d}) \\
  &= c + 0 = c.
\end{align*}
Conversely any two solutions differ by a vector in the nullspace of $(a,b)$, which forces the stated parameterization.
\end{proof}

\begin{example}
  Solve
  \[14x + 21y = 7.\]
  Here $d=\gcd(14,21)=7$ and $7\mid7$.  Extended Euclid yields
  \[14(-1)+21(1)=7,\]
  so one solution is $(-1,1)$.  Then the general solution is
  \[
  x=-1+\frac{21}{7}t = -1+3t,
  \quad
  y=1-\frac{14}{7}t = 1-2t,
  \quad
  t\in\mathbb Z.
  \]
\end{example}